\documentclass[11pt]{article}
\usepackage[margin=1in]{geometry}
\usepackage{amsmath,amssymb,amsthm}
\usepackage{booktabs,longtable,array}
\usepackage[hidelinks]{hyperref}

\newtheorem{definition}{Definition}
\newtheorem{theorem}{Theorem}
\newtheorem{proposition}{Proposition}
\newtheorem{corollary}{Corollary}
\newtheorem{example}{Example}
\newcommand{\powerset}{\mathcal{P}}

\title{Scoped Negative-Result Success and Minimum-Extension Pathway\\
\large P14-T05 draft/control formalization}
\author{The \AE ther-Flow Research Project}
\date{2026-07-22}

\begin{document}
\maketitle

\begin{abstract}
This task-local packet defines when a negative scientific result may complete
an exact research question without becoming a global rejection of the
\AE ther-Flow programme.  It distinguishes finite countermodels, scoped
obstructions, scoped no-go theorems, minimum-extension theorems, and the much
stronger claim that no future extension could succeed.  A finite hitting-set
theorem converts an explicit family of obstruction clauses into
inclusion-minimal extension signatures.  The theorem is a lower-bound and
routing result, not an ontology proposal, adoption, physical interpretation,
or proof of general source equivalence.  The current EqSrc selector and
quotient results are classified conservatively: several exact routes are
locally frozen, while general EqSrc and same-milestone source-extension
continuation remain open.
\end{abstract}

\section{Authority and exact claim boundary}

This artifact has status \texttt{draft/control}.  Its mathematical statements
are proved only in the explicit finite or scope-indexed control models defined
below.  It does not edit or supersede canonical ontology, choose a physical
gauge, adopt a source-side law, alter the Distance-to-GR ledger, establish
general EqSrc, derive a metric, couple matter, derive Einstein equations,
authorize publication, or establish a completed first-principles GR
derivation.  Exact-GR recovery remains an adopted downstream closure in the
canonical interpretation; it is not thereby a derivation from substrate
structure.

The governing discipline is:
\begin{quote}
Close the exact quantified question that the witness or theorem answers;
preserve every explicitly stated assumption and domain boundary; and name a
reopening condition.  Never promote ``this route fails'' to ``the theory
fails'' without a theorem whose quantified universe actually supports that
conclusion.
\end{quote}

\section{Scope-bearing result records}

\begin{definition}[Exact question record]
An exact question record is
\[
 Q=(I,\mathcal{M},A,C,S),
\]
where $I$ is the instance domain, $\mathcal{M}$ is the explicitly admissible
construction family, $A$ is the assumption list, $C(m,i)$ is the tested claim,
and $S$ is the immutable source-and-version fingerprint.  Two records with a
different component are different questions.
\end{definition}

\begin{definition}[Negative-result record]
A negative-result record is
\[
 R=(Q,L,W,P,O),
\]
where $L$ is a strength label from the ladder below, $W$ is a finite witness or
counterexample set when applicable, $P$ is a proof or audit record, and $O$ is
an explicit reopening predicate.  A record is \emph{scientifically complete
for its exact scope} only if all five components are present and its evidence
supports the selected strength label.
\end{definition}

The phrase ``scientifically complete for its exact scope'' is deliberately
separate from ontology adoption, physical interpretation, Gate Chair verdict,
external publication, and programme completion.

\section{Negative-result strength ladder}

\begin{definition}[Finite countermodel]
For a fixed question record $Q$, a finite countermodel is one pair
$(m,i)\in\mathcal{M}\times I$ with $\neg C(m,i)$.  It refutes a universal
claim over the stated domain.  It says nothing about constructions or
instances outside that domain.
\end{definition}

\begin{definition}[Scoped obstruction]
A scoped obstruction identifies an explicit burden $B$ and proves or audits
that a named family $\mathcal{F}\subseteq\mathcal{M}$ does not discharge $B$
under $A$.  It may freeze the exact family $\mathcal{F}$, but it is weaker than
a no-go theorem for all of $\mathcal{M}$ unless $\mathcal{F}=\mathcal{M}$ is
itself established.
\end{definition}

\begin{definition}[Scoped no-go theorem]
A scoped no-go theorem proves
\[
  \forall m\in\mathcal{M}\;\forall i\in I,\quad \neg C(m,i)
\]
or an explicitly equivalent quantified statement, under the assumptions $A$
and source fingerprint $S$.  Its conclusion is exhausted by those
quantifiers.  The theorem does not range over unlisted future extensions.
\end{definition}

\begin{definition}[Minimum-extension theorem]
A minimum-extension theorem provides necessary added-capability signatures for
escaping a finite, explicit obstruction system.  ``Minimum'' means
inclusion-minimal unless a cost function is separately declared; it does not
mean ontologically preferred, physically realized, or adopted.
\end{definition}

\begin{definition}[Future-extension impossibility]
A future-extension impossibility claim quantifies over a justified complete
universe of admissible future source extensions and proves that none can meet
the target.  This strength requires a completeness theorem for that universe
or an equivalent meta-theorem.  No such theorem is supplied here, so this
label is blocked for the present EqSrc frontier.
\end{definition}

The ladder is not a rhetoric scale.  Each step changes the quantified domain
or conclusion and therefore demands new evidence.

\section{Minimum-extension hitting-set theorem}

Let $U$ be a finite declared universe of extension capabilities.  Let
$J=\{1,\ldots,n\}$ index obstruction burdens, and let each nonempty
$C_j\subseteq U$ contain the capabilities that can discharge burden $j$.
An extension signature is a set $E\subseteq U$.

\begin{definition}[Escape predicate]
The task-local escape predicate is
\[
 \operatorname{Esc}(E)\;\Longleftrightarrow\;
 \bigwedge_{j\in J}(E\cap C_j\neq\varnothing).
\]
Thus an escaping signature must supply at least one declared capability for
each obstruction clause.  This definition states necessary control burdens;
it does not assert that satisfying them is sufficient for physics.
\end{definition}

\begin{theorem}[Finite minimum-extension signatures]
\label{thm:hitting}
For finite $U$ and nonempty clauses $C_j$, the following hold.
\begin{enumerate}
  \item $\operatorname{Esc}(E)$ holds exactly when $E$ is a hitting set of
  $\{C_j:j\in J\}$.
  \item Every escaping $E$ contains an inclusion-minimal escaping subset
  $H\subseteq E$.
  \item The collection
  \[
    \mathcal{H}_{\min}=
    \{H\subseteq U:\operatorname{Esc}(H)\ \text{and no proper subset of }H
      \text{ escapes}\}
  \]
  is finite and nonempty.
  \item Consequently, any construction that escapes all declared clauses must
  realize at least one signature $H\in\mathcal{H}_{\min}$, provided its added
  capabilities are represented in $U$ and the clause model is sound for the
  exact question record.
\end{enumerate}
\end{theorem}

\begin{proof}
Part 1 is the definition of a hitting set.  For Part 2, begin with finite $E$.
If a proper subset still escapes, delete a dispensable element; repeat.  Since
$E$ is finite, deletion terminates at an escaping subset with no escaping
proper subset.  Part 3 follows because $\powerset(U)$ is finite, $U$ itself
hits every nonempty $C_j$, and Part 2 supplies at least one minimal member.
For Part 4, represent a successful construction by its capability set $E$.
Soundness of the clauses gives $\operatorname{Esc}(E)$; Part 2 then supplies
$H\subseteq E$.  The representation and soundness qualifications are
essential and prevent an unlisted extension from being ruled out.
\end{proof}

\begin{corollary}[Weighted routing, not adoption]
If a nonnegative cost $w:U\to\mathbb{R}_{\geq0}$ is declared, one may rank
$\mathcal{H}_{\min}$ by $\sum_{u\in H}w(u)$.  A least-cost signature is least
only in that declared bookkeeping model.  It is not thereby the true ontology,
a physical law, or an authorized implementation choice.
\end{corollary}

\begin{proposition}[No global conclusion from an incomplete universe]
If a capability $u_\star$ or a sound discharge mechanism is omitted from $U$
or from every relevant clause, Theorem~\ref{thm:hitting} provides no verdict
about extensions using $u_\star$.  Therefore the theorem cannot establish
future-extension impossibility without a separate completeness result.
\end{proposition}

\begin{proof}
The theorem quantifies only over subsets of $U$ and assumes the clauses are
sound.  An omitted capability is outside that quantified universe, so no
conclusion about it follows.
\end{proof}

\section{Finite control examples}

Let $U=\{d,n,q\}$ with clauses $C_1=\{d,q\}$ and
$C_2=\{n,q\}$.  Here $d$ denotes a source-derived discriminator, $n$ a
naturality-and-variation law, and $q$ a quotient plus interface-descent
capability.  Then
\[
 \mathcal{H}_{\min}=\bigl\{\{q\},\{d,n\}\bigr\}.
\]
The empty signature and the singletons $\{d\}$ and $\{n\}$ fail.  The
signature $\{q\}$ is minimal in this toy clause model but is not automatically
physically adequate: an operational bridge and exact-GR recovery obligations
remain separate.  Likewise $\{d,n\}$ merely records necessary declared
capabilities and does not select or adopt their laws.

This two-clause example separates four claims:
\begin{enumerate}
  \item a current family can fail a burden;
  \item the exact family can be locally frozen;
  \item an explicit capability model can identify minimal reopening routes;
  \item no conclusion follows about every future extension.
\end{enumerate}

\section{Conservative classification of the EqSrc frontier}

The current records support the following bounded classification.

\begin{longtable}{>{\raggedright\arraybackslash}p{0.23\textwidth}
                  >{\raggedright\arraybackslash}p{0.24\textwidth}
                  >{\raggedright\arraybackslash}p{0.43\textwidth}}
\toprule
Existing result & Allowed strength & Preserved boundary \\
\midrule
\endhead
Empty fixed locus for an exact deterministic equivariant selector domain &
Scoped no-selector conclusion for that domain &
No conclusion about stochastic selectors, extra marks, altered morphisms,
changed variations, or other source extensions. \\
Multiple fixed loci or distinct $K$-images &
Finite countermodel or scoped nonuniqueness obstruction &
No unique physical representative and no physical-gauge verdict follow. \\
Singleton product of fixed loci &
Conditional existence and uniqueness within the declared selector model &
Existence is not source derivation, naturality, operational meaning, or
physical realization. \\
Five historical selector families and seven exact candidates &
Local family freeze under recorded assumptions &
Their coordinates are preserved; general EqSrc remains undischarged. \\
Generic invariant-functor quotient candidate &
Conditional theorem plus locally frozen generic adequacy route &
No physical quotient, interface descent, operational completeness, or
general EqSrc is established. \\
P14-T03 no-target result &
Task-local claim-logic theorem &
No-target purity is necessary and insufficient; it is not a theorem about
nature or a protected gate verdict. \\
\bottomrule
\end{longtable}

Internal AI review of the P2 selector theorem found no blocking defect in its
declared proposal-only domain.  That review is not independent external human
review and does not supply theorem, adoption, or publication authority.

\section{Provisional EqSrc extension signatures}

The following signatures are routing aids.  They are non-exhaustive,
proposal-only, and conditional on the present burden decomposition.

\begin{description}
  \item[Selection path.] A source-derived discriminator together with explicit
  morphism, transformation, and variation laws.  The discriminator must not be
  defined by importing the target metric atlas it is meant to derive.
  \item[Representative-irrelevance path.] A physical-observable functor plus a
  theorem that the admitted physical morphisms preserve and reflect the
  relevant operational content.  Algebraic orbit constancy alone is not an
  operational theorem.
  \item[Quotient path.] A source-defined invariant functor, interface descent,
  and a physical bridge that establishes why the quotient carries the required
  observables.  A mathematically defined quotient alone does not establish
  physical adequacy.
\end{description}

Stochastic selection, added marks, changed admissible variations, and altered
morphism classes are distinct extensions.  They reopen the exact route only
when their new assumptions and capability costs are declared.  They are not
counterexamples to a theorem whose scope excluded them.

\section{Closure and reopening policy}

An exact route may be marked \texttt{close\_exact\_route} when its negative
record identifies the full question tuple, evidence, strength, and reopening
predicate.  A named finite family may be marked \texttt{freeze\_family} when
every listed member has a qualifying result and the family inventory is
sealed.  A scoped question may be marked \texttt{close\_scoped\_question}
only when a theorem covers the declared admissible class.  None of these
statuses closes the entire research programme.

Reopening is evidence-triggered, not rhetorical.  Valid triggers include a
changed assumption, enlarged instance domain, new admissible construction,
new capability outside the prior universe, repaired proof defect, stronger
operational bridge, or independent review that identifies a material gap.
Every reopened route receives a new exact question record and retains the old
negative result as historical evidence.

\section{Publication-readiness criteria}

A negative result is internally manuscript-ready only if:
\begin{enumerate}
  \item the exact question tuple and immutable source hashes are included;
  \item the selected strength label is supported by a checkable proof,
  countermodel, or sealed audit;
  \item scope, assumptions, non-implications, and reopening criteria are stated;
  \item target-atlas leakage and circularity have been audited where relevant;
  \item mathematical payload is separated from validator and workflow receipts;
  \item ontology, physical, ledger, review, and publication authority are named
  separately;
  \item citations and related negative results are preserved; and
  \item a human publication owner independently authorizes submission after
  any required external review.
\end{enumerate}

Items 1--7 are evidence-readiness criteria.  Item 8 is protected publication
authority.  This packet satisfies neither external-review nor human-submission
authority and therefore records only \texttt{draft/control} status.  ``Ready
for internal manuscript integration'' must never be shortened to ``published''
or ``publication authorized.''

\section{Distance-to-GR disposition}

The P14-T05 burden is discharged at task-local policy scope: negative science
can close an exact line while same-milestone continuation stays open.  The
scientific frontier remains unchanged.  Source ontology is not derived;
general EqSrc is undischarged; effective metric, matter coupling, Einstein
equations, and benchmark promotion remain blocked by their existing burdens.
No row of the Distance-to-GR ledger is changed.

\section{Source basis}

The controlling sources are the v21 implementation plan (P14-T05), the
canonical foundations and geometry manuscripts, the epistemic-category and
scoped-positive vocabularies, the GR derivation burden map, the Distance-to-GR
ledger, P2-T08 review evidence, the P3-T07 family-freeze completion and repair,
and the P14-T03 no-target sufficiency policy.  Their hashes and exact paths are
recorded in the associated AgentJob and scope/reopening schema.  Generated
wiki notes and local retrieval caches are derivative discovery aids only.

\end{document}
